% GENERATED FILE. Edit canonical lessons, then run npm run generate:books.
% canonical-source-hash: fnv1a64:8564529757b73e48
% canonical-lessons: LA-C362-refello, LA-C362-refuto, LA-C362-redarguo, LA-C362-convinco, LA-C362-coniectura

\chapter{Proving Someone Wrong}
\label{ch:la-proving-someone-wrong}
\hlchaptermodality{362}

\begin{chapteropening}
\textbf{By the end of this chapter:} \emph{I can refute a claim, prove someone wrong and say that a case rests only on a guess.}

Complete the last lesson of chapter 362: I can refute a claim, prove someone wrong and say that a case rests only on a guess.
\end{chapteropening}

\section[refellō, refellere]{refellō, refellere — to prove false, to refute}

\label{lesson:LA-C362-refello}



\begin{quote}
Before the new one: say the Latin for to show, then the Latin for to strengthen.
\end{quote}

\subsection*{You'll want to know: refellō, refellere}
\textbf{refellō, refellere} — \textquotedblleft{}to prove false, to refute\textquotedblright{}.

\begin{quote}
Facile erat eius verba refellere.
\end{quote}

\begin{quote}
\emph{It was easy to refute his words.}
\end{quote}

\emph{refellō, refellere, refellī}, with no supine. Built on \emph{fallō}, to deceive: to show up a deception.

\begin{cousinweb}
\emph{refellō} is built from \emph{re-} \textquotedblleft{}back\textquotedblright{} and \emph{fallō} \textquotedblleft{}to deceive\textquotedblright{} (EtymDB; Lewis \& Short): to turn a deception back on itself.
\end{cousinweb}

\begin{grammarlens}[title={What you've built}]
One more word for telling, explaining and describing at B1.
\end{grammarlens}

\subsection*{Your turn}
\begin{itemize}
\raggedright
  \item \emph{Say it:} \emph{refellō, refellere}
  \item \emph{Say it:} \emph{refellō, refellere}, once more
  \item \emph{Say it:} a sentence of your own with \emph{refellō, refellere}
  \item \emph{Recall:} say the Latin for to show, then the Latin for to strengthen, then say \emph{refellō, refellere} again
\end{itemize}

\begin{tcolorbox}[breakable,colback=teal!4,colframe=teal!35!black,title={Before you move on}]
What does refellō, refellere mean? (\textquotedblleft{}To prove false, to refute\textquotedblright{}.) Say it once more.
\end{tcolorbox}
\section[refūtō, refūtāre]{refūtō, refūtāre — to refute; to drive back}

\label{lesson:LA-C362-refuto}



\begin{quote}
Before the new one: say the Latin for to strengthen, then the Latin for to prove false.
\end{quote}

\subsection*{You'll want to know: refūtō, refūtāre}
\textbf{refūtō, refūtāre} — \textquotedblleft{}to refute; to drive back\textquotedblright{}.

\begin{quote}
Hanc opīniōnem nōn difficile est refūtāre.
\end{quote}

\begin{quote}
\emph{This opinion is not hard to refute.}
\end{quote}

\emph{refūtō, refūtāre, refūtāvī, refūtātum}. First to beat back an attack, then an argument. English \emph{refute}.

\begin{cousinweb}
\emph{refūtō} is built from \emph{re-} \textquotedblleft{}back\textquotedblright{} and an old verb \emph{fūtō}, which EtymDB glosses \textquotedblleft{}to beat\textquotedblright{}; Lewis \& Short connect it instead with \emph{fundō} \textquotedblleft{}to pour\textquotedblright{}, so its deeper origin is not settled. English \emph{refute} comes from it.
\end{cousinweb}

\begin{grammarlens}[title={What you've built}]
One more word for telling, explaining and describing at B1.
\end{grammarlens}

\subsection*{Your turn}
\begin{itemize}
\raggedright
  \item \emph{Say it:} \emph{refūtō, refūtāre}
  \item \emph{Say it:} \emph{refūtō, refūtāre}, once more
  \item \emph{Say it:} a sentence of your own with \emph{refūtō, refūtāre}
  \item \emph{Recall:} say the Latin for to strengthen, then the Latin for to prove false, then say \emph{refūtō, refūtāre} again
\end{itemize}

\begin{tcolorbox}[breakable,colback=teal!4,colframe=teal!35!black,title={Before you move on}]
What does refūtō, refūtāre mean? (\textquotedblleft{}To refute; to drive back\textquotedblright{}.) Say it once more.
\end{tcolorbox}
\section[redarguō, redarguere]{redarguō, redarguere — to disprove, to show to be wrong}

\label{lesson:LA-C362-redarguo}



\begin{quote}
Before the new one: say the Latin for to prove false, then the Latin for to refute.
\end{quote}

\subsection*{You'll want to know: redarguō, redarguere}
\textbf{redarguō, redarguere} — \textquotedblleft{}to disprove, to show to be wrong\textquotedblright{}.

\begin{quote}
Testis falsum dīxit, sed iūdex eum redarguit.
\end{quote}

\begin{quote}
\emph{The witness told a lie, but the judge proved him wrong.}
\end{quote}

\emph{redarguō, redarguere, redarguī}. Built on \emph{arguō}: to prove the other side wrong, back to its face.

\begin{cousinweb}
\emph{redarguō} is built from \emph{red-} \textquotedblleft{}back\textquotedblright{} and \emph{arguō} \textquotedblleft{}to make clear, to prove\textquotedblright{} (Lewis \& Short): to prove a claim back against the one who made it.
\end{cousinweb}

\begin{grammarlens}[title={What you've built}]
One more word for telling, explaining and describing at B1.
\end{grammarlens}

\subsection*{Your turn}
\begin{itemize}
\raggedright
  \item \emph{Say it:} \emph{redarguō, redarguere}
  \item \emph{Say it:} \emph{redarguō, redarguere}, once more
  \item \emph{Say it:} a sentence of your own with \emph{redarguō, redarguere}
  \item \emph{Recall:} say the Latin for to prove false, then the Latin for to refute, then say \emph{redarguō, redarguere} again
\end{itemize}

\begin{tcolorbox}[breakable,colback=teal!4,colframe=teal!35!black,title={Before you move on}]
What does redarguō, redarguere mean? (\textquotedblleft{}To disprove, to show to be wrong\textquotedblright{}.) Say it once more.
\end{tcolorbox}
\section[convincō, convincere]{convincō, convincere — to prove wrong, to convict; to prove}

\label{lesson:LA-C362-convinco}



\begin{quote}
Before the new one: say the Latin for to refute, then the Latin for to disprove.
\end{quote}

\subsection*{You'll want to know: convincō, convincere}
\textbf{convincō, convincere} — \textquotedblleft{}to prove wrong, to convict; to prove\textquotedblright{}.

\begin{quote}
Eum mendāciī convīcērunt.
\end{quote}

\begin{quote}
\emph{They convicted him of lying.}
\end{quote}

\emph{convincō, convincere, convīcī, convictum}. The charge goes in the genitive: \emph{mendāciī convincere}, to convict of lying. In Latin it proves, it does not persuade. English \emph{convict} and \emph{convince}.

\begin{cousinweb}
\emph{convincō} is built from \emph{con-} and \emph{vincō} \textquotedblleft{}to conquer\textquotedblright{} (EtymDB; Lewis \& Short): to beat an opponent completely in argument. English \emph{convict} and \emph{convince} both come from it.
\end{cousinweb}

\begin{grammarlens}[title={What you've built}]
One more word for telling, explaining and describing at B1.
\end{grammarlens}

\subsection*{Your turn}
\begin{itemize}
\raggedright
  \item \emph{Say it:} \emph{convincō, convincere}
  \item \emph{Say it:} \emph{convincō, convincere}, once more
  \item \emph{Say it:} a sentence of your own with \emph{convincō, convincere}
  \item \emph{Recall:} say the Latin for to refute, then the Latin for to disprove, then say \emph{convincō, convincere} again
\end{itemize}

\begin{tcolorbox}[breakable,colback=teal!4,colframe=teal!35!black,title={Before you move on}]
What does convincō, convincere mean? (\textquotedblleft{}To prove wrong, to convict; to prove\textquotedblright{}.) Say it once more.
\end{tcolorbox}
\section[coniectūra, coniectūrae]{coniectūra, coniectūrae — a guess, a conjecture; an inference}

\label{lesson:LA-C362-coniectura}



\begin{quote}
Before the new one: say the Latin for to disprove, then the Latin for to prove wrong.
\end{quote}

\subsection*{You'll want to know: coniectūra, coniectūrae}
\textbf{coniectūra, coniectūrae} — \textquotedblleft{}a guess, a conjecture; an inference\textquotedblright{}.

\begin{quote}
Hoc nōn sciō, sed coniectūrā tantum dīcō.
\end{quote}

\begin{quote}
\emph{I do not know this; I say it only by guesswork.}
\end{quote}

\emph{coniectūra, coniectūrae} f. Built on \emph{coniciō}, to throw together: pieces thrown together to make a guess. English \emph{conjecture}.

\begin{cousinweb}
\emph{coniectūra} is built on \emph{coniciō} \textquotedblleft{}to throw together, to infer\textquotedblright{} (EtymDB; Lewis \& Short), from \emph{con-} and \emph{iaciō} \textquotedblleft{}to throw\textquotedblright{}. English \emph{conjecture} comes from it.
\end{cousinweb}

\begin{grammarlens}[title={What you've built}]
One more word for telling, explaining and describing at B1.
\end{grammarlens}

\subsection*{Your turn}
\begin{itemize}
\raggedright
  \item \emph{Say it:} \emph{coniectūra, coniectūrae}
  \item \emph{Say it:} \emph{coniectūra, coniectūrae}, once more
  \item \emph{Say it:} a sentence of your own with \emph{coniectūra, coniectūrae}
  \item \emph{Recall:} say the Latin for to disprove, then the Latin for to prove wrong, then say \emph{coniectūra, coniectūrae} again
\end{itemize}

\begin{tcolorbox}[breakable,colback=teal!4,colframe=teal!35!black,title={Before you move on}]
What does coniectūra, coniectūrae mean? (\textquotedblleft{}A guess, a conjecture; an inference\textquotedblright{}.) Say it once more.
\end{tcolorbox}
